\documentclass{requirements-lab}

\hypersetup{
    pdftitle={HLA Requirements Lab Design System v0.1 — Style Specimen},
    pdfauthor={HLA Requirements Lab},
    pdfsubject={Shared publishing components for HLA Requirements Lab publications},
    pdfkeywords={Requirements Lab, design system, typography, tables, HLA}
}

\renewcommand{\RLPublicationEdition}{Design System v0.1}
\renewcommand{\RLPublicationNumber}{1516.2 family}
\renewcommand{\RLPublicationTitle}{Style specimen and component gallery}
\renewcommand{\RLPublicationDescriptor}{Shared technical publishing system for HLA Requirements Lab editions}
\renewcommand{\RLPublicationSource}{the validated HLA 1516-family semantic corpus}
\renewcommand{\RLPublicationStatus}{Prototype presentation layer; no normative content is defined here.}
\renewcommand{\RLClauseContext}{STYLE SPECIMEN}

\begin{document}

\RLCoverPage
\RLDocumentStatusPage
\RLHowToReadPage

\tableofcontents
\clearpage

\section{Design principles}
\RLAccentRule
\RLBodyFont
The Requirements Lab publishing system separates semantic ownership from visual presentation. A semantic record says whether something is a clause, requirement, table, note, example, figure, or reference; the house style decides how that record appears in print and on screen.

The system is intentionally restrained. It should feel like a contemporary systems-engineering reference, language specification, and standards handbook—not a corporate report and not a facsimile of an external publisher's document.

\RLNote{A note is informative content. Its treatment provides a visual cue without changing the status of the surrounding prose.}

\RLExample{An example illustrates a possible representation. It is not a new requirement.}

\RLLabNote{This specimen is Lab-authored design-system documentation. It is not source-derived technical content.}

\clearpage
\section{Typography and technical tokens}
\RLAccentRule
\RLBodyFont
Narrative prose remains calm and readable. Navigation and headings use a sans-serif voice.

Machine-significant tokens use a compact monospace voice so the reader can distinguish the language of the specification from the language describing it.

\subsection{Token examples}
The root class is \RLClass{HLAobjectRoot}. An attribute such as \RLAttribute{HLAprivilegeToDeleteObject} is distinct from the prose around it. A service such as \RLService{Update Attribute Values} remains readable, while an XML element such as \RLXML{objectClass} and a value such as \RLValue{TimeStamp} remain visibly machine-oriented.

\subsection{Body and table scale}
\RLBodyFont Body prose uses approximately 9.75–10 points with generous leading. \par
\RLTableFont Table text uses approximately 8–9 points. \par
\RLSmallFont Notes and captions use approximately 8–8.5 points.

\clearpage
\section{Visible hierarchy}
\RLAccentRule
\RLBodyFont
The hierarchy is structural first and visual second. Major clauses receive a strong opener; subclauses retain their numbers but use less visual weight as depth increases.

\subsection{Major subclause}
\subsubsection{Nested heading}
\paragraph{Deep heading} The descriptive title remains readable while the number becomes secondary.

\RLClauseOpener{04}{HLA OMT COMPONENTS}{Object models consist of interrelated components describing classes, attributes, interactions, parameters, and representations.}

\section{Notes, examples, and editorial material}
\RLAccentRule
\RLBodyFont
Source-derived material and Lab-authored interpretation should never be blended invisibly. The callouts below use the same family of colors and typography while declaring their role.

\RLNote{Normative status belongs to the semantic content model. A visual note treatment does not, by itself, make surrounding content informative.}

\RLExample{A reader can scan the page and locate an illustration without mistaking it for a conformance obligation.}

\RLRationale{A rationale can explain a presentation or design choice where the Lab explicitly authors that explanation.}

\RLLabNote{The Lab note is the only callout in this specimen that is explicitly editorial. It is kept visually separate from source-derived notes.}

\clearpage
\section{Reference tables}
\RLAccentRule
\RLBodyFont
Simple reference tables favor whitespace and horizontal grouping over a heavy spreadsheet grid.

\begin{RLReferenceTable}{Simple reference table}{tab:reference}
    \RLIdentifier{HLAobjectRoot} & Root object class from which object classes are derived. \\
    \RLIdentifier{HLAinteractionRoot} & Root interaction class from which interaction classes are derived. \\
    \RLIdentifier{HLAinteger64Time} & Example of a machine-significant time representation. \\
\end{RLReferenceTable}

\begin{RLRegistryTable}{Registry/reference table}{tab:registry}
    \RLIdentifier{HLAreliable} & transport & Reliable delivery semantics. \\
    \RLIdentifier{HLAbestEffort} & transport & Best-effort delivery semantics. \\
    \RLIdentifier{TimeStamp} & value & A sample machine-significant value. \\
\end{RLRegistryTable}

\clearpage
\section{Dense normative matrix}
\RLAccentRule
\RLBodyFont
Dense matrices require stronger structure because columns have formal meaning. The specimen uses grouped horizontal rules and a deliberate identifier column rather than a black box around every cell.

\begin{center}
\begin{tabularx}{\linewidth}{@{}L{0.20\linewidth}L{0.18\linewidth}L{0.17\linewidth}Y@{}}
\toprule
\RLTableHeader{Identifier} & \RLTableHeader{Type} & \RLTableHeader{Cardinality} & \RLTableHeader{Meaning} \\
\midrule
    \RLIdentifier{HLAdeletePrivilege} & attribute & 0..1 & Indicates whether a joined federate has the privilege to delete an object instance. \\
\RLIdentifier{HLAinteractionRoot} & class & 1 & Root of the interaction-class hierarchy. \\
\RLIdentifier{HLAinteger64BE} & datatype & 1 & Big-endian signed integer representation. \\
\bottomrule
\end{tabularx}
\end{center}

\RLExample{When a matrix has more semantic columns than portrait measure can support, landscape is the correct presentation—not a layout failure.}

\clearpage
\begin{landscape}
\section{Landscape and continued tables}
\RLAccentRule
\RLBodyFont
Landscape is reserved for information whose horizontal structure carries meaning. The following table is deliberately wide enough to preserve identifier and description readability.

\begin{center}
\begin{longtable}{@{}L{0.18\linewidth}L{0.13\linewidth}L{0.12\linewidth}L{0.13\linewidth}L{0.13\linewidth}L{0.18\linewidth}@{}}
\caption{Landscape dense matrix with repeated semantic groups}\label{tab:landscape}\\
\toprule
\RLTableHeader{Identifier} & \RLTableHeader{Kind} & \RLTableHeader{Direction} & \RLTableHeader{Order} & \RLTableHeader{Cardinality} & \RLTableHeader{Description} \\
\midrule
\endfirsthead
\caption[]{Landscape dense matrix (continued)}\\
\toprule
\RLTableHeader{Identifier} & \RLTableHeader{Kind} & \RLTableHeader{Direction} & \RLTableHeader{Order} & \RLTableHeader{Cardinality} & \RLTableHeader{Description} \\
\midrule
\endhead
\RLIdentifier{UpdateAttributeValues} & service & federate $\rightarrow$ RTI & timestamp & 0..many & Sends updated attribute values for one or more object instances. \\
    \RLIdentifier{SendInteraction} & service & federate $\rightarrow$ RTI & receive & 0..many & Sends an interaction and its associated parameters. \\
\RLIdentifier{ReflectAttributeValues} & callback & RTI $\rightarrow$ federate & timestamp & 0..many & Delivers reflected attribute values to a subscribing federate. \\
    \RLIdentifier{ReceiveInteraction} & callback & RTI $\rightarrow$ federate & receive & 0..many & Delivers an interaction and its parameters. \\
\RLIdentifier{HLAinteger64Time} & datatype & value & n/a & 1 & Represents a logical time value. \\
\RLIdentifier{HLAfloat64BE} & datatype & value & n/a & 1 & Represents a binary floating-point value. \\
\bottomrule
\end{longtable}
\end{center}
\end{landscape}

\clearpage
\section{Hierarchy views}
\RLAccentRule
\RLBodyFont
Object and interaction class structures are not ordinary spreadsheets. A compact hierarchy view can expose parentage while preserving the strict tabular representation elsewhere.

\begin{tabularx}{\linewidth}{@{}L{0.34\linewidth}Y@{}}
\toprule
\RLTableHeader{Class path} & \RLTableHeader{Role} \\
\midrule
\RLClass{HLAobjectRoot} & Root object class. \\
\quad\,\RLClass{ObjectRoot} & Example child class. \\
\qquad\,\RLClass{Vehicle} & Example nested class. \\
\qquad\,\,\RLClass{Aircraft} & Example deep class. \\
\bottomrule
\end{tabularx}

\clearpage
\section{Equations and byte layouts}
\RLAccentRule
\RLBodyFont
Equations should have breathing room, aligned labels, and an explanation that can be scanned independently.

\begin{equation}
    (\RLVariable{Offset}_{i} + \RLVariable{Size}_{i} + \RLVariable{P}_{i})
    \bmod \RLVariable{V}_{i+1} = 0
    \RLEquationLabel{1}
\end{equation}

\RLDefinitionPair{Offset\textsubscript{i}}{byte offset of component \RLValue{i}, in bytes}
\RLDefinitionPair{Size\textsubscript{i}}{size of component \RLValue{i}, in bytes}
\RLDefinitionPair{V\textsubscript{i+1}}{octet-boundary value of the following component}
\RLDefinitionPair{P\textsubscript{i}}{required padding, in bytes}

\subsection{Byte-layout example}
\begin{center}
\begin{tabular}{@{}c|c|c|c|c@{}}
\toprule
\RLTableHeader{0} & \RLTableHeader{1} & \RLTableHeader{2} & \RLTableHeader{3} & \RLTableHeader{4–7} \\
\midrule
\RLValue{octet} & \RLValue{padding} & \RLValue{padding} & \RLValue{flags} & \RLType{HLAinteger32BE} \\
\bottomrule
\end{tabular}
\end{center}

\clearpage
\section{Code, XML, and identifiers}
\RLAccentRule
\RLBodyFont
Machine-oriented material should remain searchable and selectable while receiving enough visual separation from narrative prose.

\RLCodeInline{HLAobjectRoot.Vehicle.Aircraft}

\RLCodeStart
\begin{verbatim}
<objectClass name="Vehicle">
  <attribute name="HLAprivilegeToDeleteObject"
             dataType="HLAboolean"/>
</objectClass>
\end{verbatim}
\RLCodeEnd

\RLNote{The code treatment is a presentation choice. It does not introduce an implementation requirement beyond the source material represented by the example.}

\clearpage
\section{Figures and diagrams}
\RLAccentRule
\RLBodyFont
Lab-authored diagrams should have a caption, stable identifier, and enough contrast to survive grayscale printing.

\begin{figure}[htbp]
    \centering
    \begin{tikzpicture}[node distance=1.7cm, every node/.style={font=\RLSans\small}]
        \node[draw=RLBlue,fill=RLMist,rounded corners,minimum width=3.2cm,minimum height=0.8cm] (federate) {Federate};
        \node[draw=RLOrange,fill=RLMist,rounded corners,minimum width=3.2cm,minimum height=0.8cm,right=of federate] (rti) {Runtime infrastructure};
        \node[draw=RLBlue,fill=RLMist,rounded corners,minimum width=3.2cm,minimum height=0.8cm,right=of rti] (federation) {Federation};
        \draw[-{Stealth},RLBlue,line width=1pt] (federate) -- node[above]{services} (rti);
        \draw[-{Stealth},RLOrange,line width=1pt] (rti) -- node[above]{execution} (federation);
    \end{tikzpicture}
    \caption{Simple systems relationship diagram.}
    \label{fig:systems-relationship}
\end{figure}

\clearpage
\section{Cross-references and navigation}
\RLAccentRule
\RLBodyFont
The modern edition should make navigation feel like documentation. Internal links, bookmarks, clause references, and figure references should be stable and meaningful.

See \RLClauseLink{fig:systems-relationship}{Figure~\ref{fig:systems-relationship}} for the systems relationship example. A future generated edition can resolve a source identifier such as \RLIdentifier{content-15162-page-016} to a stable semantic destination.

\begin{itemize}
    \item Clickable table and figure references.
    \item Hierarchical PDF bookmarks generated from the clause tree.
    \item Meaningful PDF metadata and logical page labels.
    \item Searchable text with identifiers preserved as selectable text.
\end{itemize}

\clearpage
\section{Provenance-safe presentation}
\RLAccentRule
\RLBodyFont
The publication has three conceptual provenance classes:

\begin{description}[leftmargin=3.0cm,style=multiline]
    \item[Source-derived] Faithful reconstruction, transcription, or paraphrase from supplied source material.
    \item[Editorial apparatus] Navigation, repaired layout, source references, and formatting decisions.
    \item[Requirements Lab analysis] Semantic decomposition, quality gates, traceability products, and explicit Lab commentary.
\end{description}

\RLLabNote{The shared style makes these distinctions available without stamping every paragraph or adding legal boilerplate to every page.}

\clearpage
\section{QA contract}
\RLAccentRule
\RLBodyFont
The specimen is a visual unit test. A house-style change should be reviewed here before it is applied to a full publication.

\begin{RLRegistryTable}{Design-system QA checklist}{tab:qa}
    \RLIdentifier{RL-SEM-01} & semantics & No source record disappears during presentation migration. \\
    \RLIdentifier{RL-TBL-01} & tables & Headers repeat and cell counts remain stable. \\
    \RLIdentifier{RL-LAY-01} & layout & No clipped content, accidental blank pages, or orphan headings. \\
    \RLIdentifier{RL-NAV-01} & navigation & Bookmarks, links, metadata, and searchable text work. \\
    \RLIdentifier{RL-VIS-01} & visual & Specimen and approved pages remain visually coherent. \\
\end{RLRegistryTable}

\appendix
\section{Component inventory}
\RLAccentRule
\RLBodyFont
The v0.1 inventory is deliberately small and composable: brand tokens, type roles, page furniture, hierarchy, table taxonomy, code/XML, mathematics, notes, links, and front matter. The 1516.2 migration will add only components that are required by real source structures.

\begin{itemize}
    \item Standard profile: restrained source compilation.
    \item Analysis profile: more explicit semantic and evidence surfaces.
    \item Facsimile profile: source-oriented reconstruction for visual comparison.
\end{itemize}

\end{document}
